% answers.tex —— 2009 年·普通高等学校招生考试·北京卷·文科 · 参考答案与详细解析片段
% 由 tools/generate_projects.py 依来源总账自 2009/2009_数学真题_北京文_试题.md + 2009/2009_数学真题_北京文_解析正文.md + 2009/2009_数学真题_北京文_答案要点.md 生成，勿手改（改源 MD 后重跑）


\begin{gktip}
固定候选 PDF 仅含题面，无答案、解析和评分标准，原卷亦未标分值；以下为非官方编者独立推导，官方答案字段为空。固定源第 14 题「孤立元」定义漏印「且 $k+1\notin A$」，照录原页字面并在文中标明。
\end{gktip}

\begin{answer}
\gksolhead{1}{}

\textbf{题目}：

设集合 $A=\left\{x\mid-\dfrac12<x<2\right\}$，$B=\{x\mid x^2\leqslant1\}$，则 $A\cup B=$（　）

（A）$\{x\mid-1\leqslant x<2\}$

（B）$\left\{x\mid-\dfrac12<x\leqslant1\right\}$

（C）$\{x\mid x<2\}$

（D）$\{x\mid1\leqslant x<2\}$

\Needspace{4\baselineskip}
\textbf{解析}

\Needspace{5\baselineskip}
$B=\{x\mid x^2\leqslant1\}=[-1,1]$，$A=\left(-\frac12,2\right)$。并集为

\nopagebreak[4]
\[A\cup B=[-1,2)=\{x\mid-1\leqslant x<2\},\]

选 A。

\gkanswer{答案：A。}

\end{answer}
\begin{answer}
\gksolhead{2}{}

\textbf{题目}：

已知向量 $\boldsymbol a=(1,0)$，$\boldsymbol b=(0,1)$，$\boldsymbol c=k\boldsymbol a+\boldsymbol b$（$k\in\mathbf R$），$\boldsymbol d=\boldsymbol a-\boldsymbol b$．如果 $\boldsymbol c\parallel\boldsymbol d$，那么（　）

（A）$k=1$ 且 $\boldsymbol c$ 与 $\boldsymbol d$ 同向

（B）$k=1$ 且 $\boldsymbol c$ 与 $\boldsymbol d$ 反向

（C）$k=-1$ 且 $\boldsymbol c$ 与 $\boldsymbol d$ 同向

（D）$k=-1$ 且 $\boldsymbol c$ 与 $\boldsymbol d$ 反向

\Needspace{4\baselineskip}
\textbf{解析}

\Needspace{5\baselineskip}
$\boldsymbol c=k(1,0)+(0,1)=(k,1)$，$\boldsymbol d=(1,-1)$。由 $\boldsymbol c\parallel\boldsymbol d$：

\nopagebreak[4]
\[k\cdot(-1)-1\cdot1=0\ \Longrightarrow\ k=-1 .\]

此时 $\boldsymbol c=(-1,1)=-(1,-1)=-\boldsymbol d$，两向量反向，选 D。

\gkanswer{答案：D。}

\end{answer}
\begin{answer}
\gksolhead{3}{}

\textbf{题目}：

若 $(1+\sqrt2)^4=a+b\sqrt2$（$a,b$ 为有理数），则 $a+b=$（　）

（A）$33$

（B）$29$

（C）$23$

（D）$19$

\Needspace{4\baselineskip}
\textbf{解析}

\nopagebreak[4]
\[(1+\sqrt2)^2=1+2\sqrt2+2=3+2\sqrt2,\]

\nopagebreak[4]
\[(3+2\sqrt2)^2=9+12\sqrt2+8=17+12\sqrt2 .\]

故 $a=17$、$b=12$，$a+b=29$，选 B。

\gkanswer{答案：B。}

\end{answer}
\begin{answer}
\gksolhead{4}{}

\textbf{题目}：

为了得到函数 $y=\lg\dfrac{x+3}{10}$ 的图象，只需把函数 $y=\lg x$ 的图象上所有的点（　）

（A）向左平移 $3$ 个单位长度，再向上平移 $1$ 个单位长度

（B）向右平移 $3$ 个单位长度，再向上平移 $1$ 个单位长度

（C）向左平移 $3$ 个单位长度，再向下平移 $1$ 个单位长度

（D）向右平移 $3$ 个单位长度，再向下平移 $1$ 个单位长度

\Needspace{4\baselineskip}
\textbf{解析}

\nopagebreak[4]
\[y=\lg\frac{x+3}{10}=\lg(x+3)-1 .\]

由 $y=\lg x$ 向左平移 $3$ 得 $y=\lg(x+3)$，再向下平移 $1$ 得所求图象，选 C。

\gkanswer{答案：C。}

\end{answer}
\begin{answer}
\gksolhead{5}{}

\textbf{题目}：

用数字 $1,2,3,4,5$ 组成的无重复数字的四位偶数的个数为（　）

（A）$8$

（B）$24$

（C）$48$

（D）$120$

\Needspace{4\baselineskip}
\textbf{解析}

\Needspace{5\baselineskip}
四位偶数末位取 $2$ 或 $4$，共 $2$ 种；前三位从剩余 $4$ 个数字中选三个排列，有 $4\times3\times2=24$ 种。故共

\nopagebreak[4]
\[2\times24=48\]

个，选 C。

\gkanswer{答案：C。}

\end{answer}
\begin{answer}
\gksolhead{6}{}

\textbf{题目}：

「$\alpha=\dfrac{\pi}{6}$」是「$\cos2\alpha=\dfrac12$」的（　）

（A）充分而不必要条件

（B）必要而不充分条件

（C）充分必要条件

（D）既不充分也不必要条件

\Needspace{4\baselineskip}
\textbf{解析}

$\alpha=\frac\pi6$ 时 $\cos2\alpha=\cos\frac\pi3=\frac12$，故充分。

\Needspace{5\baselineskip}
反之，$\cos2\alpha=\frac12$ 得

\nopagebreak[4]
\[2\alpha=\pm\frac\pi3+2k\pi\ \Longrightarrow\ \alpha=\pm\frac\pi6+k\pi,\]

不一定是 $\frac\pi6$，故不必要。

所以是充分而不必要条件，选 A。

\gkanswer{答案：A。}

\end{answer}
\begin{answer}
\gksolhead{7}{}

\textbf{题目}：

若正四棱柱 $ABCD-A_1B_1C_1D_1$ 的底面边长为 $1$，$AB_1$ 与底面 $ABCD$ 成 $60^\circ$ 角，则 $A_1C_1$ 到底面 $ABCD$ 的距离为（　）

（A）$\dfrac{\sqrt3}{3}$

（B）$1$

（C）$\sqrt2$

（D）$\sqrt3$

\Needspace{4\baselineskip}
\textbf{解析}

\Needspace{5\baselineskip}
$AB_1$ 在底面的投影为 $AB=1$，与底面夹角 $60^\circ$，故高

\nopagebreak[4]
\[h=AB\tan60^\circ=\sqrt3 .\]

$A_1C_1\parallel$ 底面，其到底面的距离等于高 $\sqrt3$，选 D。

\gkanswer{答案：D。}

\end{answer}
\begin{answer}
\gksolhead{8}{}

\textbf{题目}：

设 $D$ 是正 $\triangle P_1P_2P_3$ 及其内部的点构成的集合，点 $P_0$ 是 $\triangle P_1P_2P_3$ 的中心．若集合 $S=\{P\mid P\in D,\ |PP_0|\leqslant|PP_i|,\ i=1,2,3\}$，则集合 $S$ 表示的平面区域是（　）

（A）三角形区域

（B）四边形区域

（C）五边形区域

（D）六边形区域

\Needspace{4\baselineskip}
\textbf{解析}

$S$ 是到 $P_0$ 的距离不大于到三个顶点距离的点集，即 $P_0$ 与 $P_1,P_2,P_3$ 四个点的 Voronoi 图中 $P_0$ 的胞腔与正三角形的交集。

$P_0$ 到三顶点等距，三条垂直平分线两两相交，把中心附近围成一个正六边形；该六边形完全落在正三角形内。故 $S$ 为六边形区域，选 D。

\gkanswer{答案：D。}

\end{answer}
\begin{answer}
\gksolhead{9}{}

\textbf{题目}：

若 $\sin\theta=-\dfrac45$，$\tan\theta>0$，则 $\cos\theta=$\gkblank{}．

\Needspace{4\baselineskip}
\textbf{解析}

\Needspace{5\baselineskip}
$\sin\theta<0$ 且 $\tan\theta>0$，故 $\theta$ 在第三象限，$\cos\theta<0$：

\nopagebreak[4]
\[\cos\theta=-\sqrt{1-\sin^2\theta}=-\sqrt{1-\frac{16}{25}}=-\frac35 .\]

\gkanswer{答案：$-\frac35$。}

\end{answer}
\begin{answer}
\gksolhead{10}{}

\textbf{题目}：

若数列 $\{a_n\}$ 满足：$a_1=1$，$a_{n+1}=2a_n$（$n\in\mathbf N^*$），则 $a_5=$\gkblank{}；前 $8$ 项的和 $S_8=$\gkblank{}．（用数字作答）

\Needspace{4\baselineskip}
\textbf{解析}

\Needspace{5\baselineskip}
$a_{n+1}=2a_n$，$a_1=1$，故 $a_n=2^{n-1}$。

\nopagebreak[4]
\[a_5=2^4=16,\qquad S_8=\frac{1-2^8}{1-2}=255 .\]

\gkanswer{答案：$16$；$255$。}

\end{answer}
\begin{answer}
\gksolhead{11}{}

\textbf{题目}：

若实数 $x,y$ 满足 $\begin{cases}x+y-2\geqslant0,\\ x\leqslant4,\\ y\leqslant5,\end{cases}$ 则 $s=x+y$ 的最大值为\gkblank{}．

\Needspace{4\baselineskip}
\textbf{解析}

\Needspace{5\baselineskip}
约束为 $y\geqslant2-x$、$x\leqslant4$、$y\leqslant5$。$s=x+y$ 随 $x,y$ 增大而增大，在 $x=4$、$y=5$ 处取得最大值（满足 $4+5-2>0$）：

\nopagebreak[4]
\[s_{\max}=9 .\]

\gkanswer{答案：$9$。}

\end{answer}
\begin{answer}
\gksolhead{12}{}

\textbf{题目}：

已知函数 $f(x)=\begin{cases}3^x,&x\leqslant1,\\ -x,&x>1.\end{cases}$ 若 $f(x)=2$，则 $x=$\gkblank{}．

\Needspace{4\baselineskip}
\textbf{解析}

若 $x\leqslant1$：$3^x=2$，$x=\log_3 2\approx0.631\leqslant1$，符合。

若 $x>1$：$-x=2$，$x=-2$，不符合 $x>1$。

故 $x=\log_3 2$。

\gkanswer{答案：$\log_3 2$。}

\end{answer}
\begin{answer}
\gksolhead{13}{}

\textbf{题目}：

椭圆 $\dfrac{x^2}{9}+\dfrac{y^2}{2}=1$ 的焦点为 $F_1,F_2$，点 $P$ 在椭圆上．若 $|PF_1|=4$，则 $|PF_2|=$\gkblank{}；$\angle F_1PF_2$ 的大小为\gkblank{}．

\Needspace{4\baselineskip}
\textbf{解析}

椭圆 $a=3$，由定义 $|PF_1|+|PF_2|=6$，故 $|PF_2|=2$。

\Needspace{5\baselineskip}
又 $c=\sqrt{9-2}=\sqrt7$，$|F_1F_2|=2\sqrt7$。由余弦定理：

\nopagebreak[4]
\[\cos\angle F_1PF_2=\frac{4^2+2^2-(2\sqrt7)^2}{2\cdot4\cdot2}
=\frac{20-28}{16}=-\frac12 .\]

故 $\angle F_1PF_2=\frac{2\pi}{3}$。

\gkanswer{答案：$2$；$\frac{2\pi}{3}$。}

\end{answer}
\begin{answer}
\gksolhead{14}{}

\textbf{题目}：

设 $A$ 是整数集的一个非空子集．对于 $k\in A$，如果 $k-1\notin A$，那么 $k$ 是 $A$ 的一个「孤立元」．给定 $S=\{1,2,3,4,5,6,7,8\}$，由 $S$ 的 $3$ 个元素构成的所有集合中，不含「孤立元」的集合共有\gkblank{}个．

\Needspace{4\baselineskip}
\textbf{解析}

按原页字面，$k$ 是孤立元只需 $k-1\notin A$。

\Needspace{5\baselineskip}
对任意有限非空集合 $A$，设 $k_{\min}$ 为最小元，则 $k_{\min}-1<k_{\min}$，故 $k_{\min}-1\notin A$，$k_{\min}$ 必为孤立元。因此任何 3 元子集都含孤立元，不含孤立元的集合有

\nopagebreak[4]
\[0\]

个。

\begin{gktip}
源缺陷：公开同题版本的定义为「$k-1\notin A$ 且 $k+1\notin A$」，原页漏印后一半。按官方定义，3 元集合不含孤立元等价于三个数连续，共 $\{1,2,3\},\cdots,\{6,7,8\}$ 六个。本库照录原页，登记 ISSUE，不擅自补写。
\end{gktip}

\gkanswer{答案：0（按原页字面）；官方同题为 6。}

\end{answer}
\begin{answer}
\gksolhead{15}{}

\textbf{题目}：

已知函数 $f(x)=2\sin(\pi-x)\cos x$．

（1）求 $f(x)$ 的最小正周期；

（2）求 $f(x)$ 在区间 $\left[-\dfrac{\pi}{6},\dfrac{\pi}{2}\right]$ 上的最大值和最小值．

\Needspace{4\baselineskip}
\textbf{解析}

\Needspace{5\baselineskip}
（1）

\nopagebreak[4]
\[f(x)=2\sin(\pi-x)\cos x=2\sin x\cos x=\sin2x .\]

\Needspace{5\baselineskip}
最小正周期

\nopagebreak[4]
\[T=\frac{2\pi}{2}=\pi .\]

\Needspace{5\baselineskip}
（2）$x\in\left[-\frac\pi6,\frac\pi2\right]$ 时，

\nopagebreak[4]
\[2x\in\left[-\frac\pi3,\pi\right],\qquad \sin2x\in\left[-\frac{\sqrt3}{2},1\right].\]

最大值 $1$ 在 $2x=\frac\pi2$（$x=\frac\pi4$）取得；最小值 $-\frac{\sqrt3}{2}$ 在 $2x=-\frac\pi3$（$x=-\frac\pi6$）取得。

\gkanswer{答案：（1）$\pi$；（2）最大值 $1$，最小值 $-\frac{\sqrt3}{2}$。}

\end{answer}
\begin{answer}
\gksolhead{16}{}

\textbf{题目}：

如图，四棱锥 $P-ABCD$ 的底面是正方形，$PD\perp ABCD$，点 $E$ 在棱 $PB$ 上．
\nopagebreak[4]
\begin{center}% 2009 年北京卷文科第 16 题：四棱锥 P-ABCD，底面正方形，PD⊥底面，E 在棱 PB 上。
% 来源 PDF：普通高考/2009/2009北京文.pdf，SHA256 5b310eb0f16d1bd48553d20e7f16b10784e6c83787ef7dd8fa96e3065ef4c159
% 原页：第 1 页（A3 三栏）第 3 栏上部，印刷页 501。
% 重绘依据：底面正方形 A 左下、B 右下、C 右上、D 左上（斜投影），P 在 D 的正上方；
%   可见棱 PA、PB、PC、AB、BC 实线，被遮棱 PD、AD、DC 虚线；E 取 PB 中点，AE、CE 细线。
\begin{tikzpicture}[scale=0.95, semithick, every node/.style={inner sep=1.5pt}]
  \coordinate (A) at (0,0);
  \coordinate (B) at (2.4,0);
  \coordinate (C) at (3.2,1.1);
  \coordinate (D) at (0.8,1.1);
  \coordinate (P) at (0.8,3.1);
  % 被遮棱
  \draw[dashed] (P) -- (D);
  \draw[dashed] (A) -- (D);
  \draw[dashed] (D) -- (C);
  \draw[dashed] (A) -- (C);
  % 可见棱
  \draw (P) -- (A);
  \draw (P) -- (B);
  \draw (P) -- (C);
  \draw (A) -- (B) -- (C);
  % E 与 AE、CE
  \coordinate (E) at ($(P)!0.5!(B)$);
  \fill (E) circle (1.3pt);
  \draw (A) -- (E);
  \draw (C) -- (E);
  % 标注
  \node[anchor=north east] at (A) {$A$};
  \node[anchor=north west] at (B) {$B$};
  \node[anchor=west] at (C) {$C$};
  \node[anchor=east] at (D) {$D$};
  \node[anchor=south] at (P) {$P$};
  \node[anchor=west] at (E) {$E$};
\end{tikzpicture}
\end{center}
\nopagebreak[4]
（1）求证：平面 $AEC\perp$ 平面 $PDB$；

（2）当 $PD=\sqrt2AB$ 且 $E$ 为 $PB$ 的中点时，求 $AE$ 与平面 $PDB$ 所成的角的大小．

\Needspace{4\baselineskip}
\textbf{解析}

\Needspace{5\baselineskip}
（1）底面 $ABCD$ 为正方形，故 $AC\perp BD$。又 $PD\perp$ 底面，故 $PD\perp AC$。于是 $AC\perp$ 平面 $PDB$。因 $AC\subset$ 平面 $AEC$，所以

\nopagebreak[4]
\[\text{\text{平面} }AEC\perp\text{\text{平面} }PDB .\]

\Needspace{5\baselineskip}
（2）设 $AB=1$，建立坐标系：

\nopagebreak[4]
\[D(0,0,0),\ A(1,0,0),\ B(1,1,0),\ C(0,1,0),\ P(0,0,\sqrt2).\]

\Needspace{5\baselineskip}
$E$ 为 $PB$ 中点：

\nopagebreak[4]
\[E\left(\frac12,\frac12,\frac{\sqrt2}{2}\right),\qquad
\overrightarrow{AE}=\left(-\frac12,\frac12,\frac{\sqrt2}{2}\right).\]

\Needspace{5\baselineskip}
平面 $PDB$ 由 $PD$（$z$ 轴）与 $DB=(1,1,0)$ 张成，法向量

\nopagebreak[4]
\[\mathbf n=(1,-1,0).\]

\Needspace{5\baselineskip}
设 $AE$ 与平面 $PDB$ 所成角为 $\theta$：

\nopagebreak[4]
\[\sin\theta=\frac{|\overrightarrow{AE}\cdot\mathbf n|}{|\overrightarrow{AE}||\mathbf n|}
=\frac{\left|-\frac12-\frac12\right|}
{\sqrt{\frac14+\frac14+\frac12}\cdot\sqrt2}
=\frac{1}{\sqrt2}=\frac{\sqrt2}{2}.\]

故 $\theta=45^\circ$。

\gkanswer{答案：（1）证明见解析；（2）$45^\circ$。}

\end{answer}
\begin{answer}
\gksolhead{17}{}

\textbf{题目}：

某学生在上学路上要经过 $4$ 个路口，假设在各路口是否遇到红灯是相互独立的，遇到红灯的概率都是 $\dfrac13$，遇到红灯时停留的时间都是 $2$ min．

（1）求这名学生在上学路上到第三个路口时首次遇到红灯的概率；

（2）这名学生在上学路上因遇到红灯停留的总时间至多是 $4$ min 的概率．

\Needspace{4\baselineskip}
\textbf{解析}

\Needspace{5\baselineskip}
（1）前两个路口未遇红灯、第三个遇红灯：

\nopagebreak[4]
\[P=\left(\frac23\right)^2\cdot\frac13=\frac4{27}.\]

\Needspace{5\baselineskip}
（2）停留总时间至多 $4$ min 等价于红灯数至多 $2$。设红灯数 $X\sim B\left(4,\frac13\right)$：

\nopagebreak[4]
\[P(X\leqslant2)=1-P(X=3)-P(X=4)
=1-\mathrm C_4^3\left(\frac13\right)^3\cdot\frac23-\left(\frac13\right)^4
=1-\frac8{81}-\frac1{81}=\frac{72}{81}=\frac89 .\]

\gkanswer{答案：（1）$\frac4{27}$；（2）$\frac89$。}

\end{answer}
\begin{answer}
\gksolhead{18}{}

\textbf{题目}：

设函数 $f(x)=x^3-3ax+b$（$a\neq0$）．

（1）若曲线 $y=f(x)$ 在点 $(2,f(2))$ 处与直线 $y=8$ 相切，求 $a,b$ 的值；

（2）求函数 $f(x)$ 的单调区间与极值点．

\Needspace{4\baselineskip}
\textbf{解析}

\Needspace{5\baselineskip}
（1）$f'(x)=3x^2-3a$。曲线在 $(2,f(2))$ 处与 $y=8$ 相切，故

\nopagebreak[4]
\[f(2)=8,\qquad f'(2)=0 .\]

\nopagebreak[4]
\[f'(2)=12-3a=0\ \Longrightarrow\ a=4,\]

\nopagebreak[4]
\[f(2)=8-24+b=8\ \Longrightarrow\ b=24 .\]

\Needspace{5\baselineskip}
（2）此时

\nopagebreak[4]
\[f(x)=x^3-12x+24,\qquad f'(x)=3(x-2)(x+2).\]

\begin{itemize}
\item $x<-2$ 或 $x>2$ 时 $f'(x)>0$，$f$ 递增；
\item $-2<x<2$ 时 $f'(x)<0$，$f$ 递减。
\end{itemize}

\Needspace{5\baselineskip}
极大值

\nopagebreak[4]
\[f(-2)=-8+24+24=40,\]

\Needspace{5\baselineskip}
极小值

\nopagebreak[4]
\[f(2)=8-24+24=8 .\]

\gkanswer{答案：（1）$a=4$，$b=24$；（2）增区间 $(-\infty,-2)$、$(2,+\infty)$，减区间 $(-2,2)$；极大值 $f(-2)=40$，极小值 $f(2)=8$。}

\end{answer}
\begin{answer}
\gksolhead{19}{}

\textbf{题目}：

已知双曲线 $C:\dfrac{x^2}{a^2}-\dfrac{y^2}{b^2}=1$（$a>0,b>0$）的离心率为 $\sqrt3$，右准线方程为 $x=\dfrac{\sqrt3}{3}$．

（1）求双曲线 $C$ 的方程；

（2）已知直线 $x-y+m=0$ 与双曲线 $C$ 交于不同的两点 $A,B$，且线段 $AB$ 的中点在圆 $x^2+y^2=5$ 上，求 $m$ 的值．

\Needspace{4\baselineskip}
\textbf{解析}

\Needspace{5\baselineskip}
（1）$e=\frac ca=\sqrt3$，即 $c=\sqrt3a$。右准线

\nopagebreak[4]
\[x=\frac{a^2}{c}=\frac{a^2}{\sqrt3a}=\frac a{\sqrt3}=\frac{\sqrt3}{3}
\ \Longrightarrow\ a=1 .\]

\Needspace{5\baselineskip}
于是 $c=\sqrt3$，$b^2=c^2-a^2=3-1=2$，双曲线为

\nopagebreak[4]
\[x^2-\frac{y^2}{2}=1 .\]

\Needspace{5\baselineskip}
（2）$y=x+m$ 代入：

\nopagebreak[4]
\[x^2-\frac{(x+m)^2}{2}=1
\ \Longrightarrow\ 2x^2-(x+m)^2=2
\ \Longrightarrow\ x^2-2mx-m^2-2=0 .\]

\Needspace{5\baselineskip}
设交点横坐标 $x_1,x_2$，则中点横坐标 $x=m$，纵坐标 $y=x+m=2m$。中点在 $x^2+y^2=5$ 上：

\nopagebreak[4]
\[m^2+(2m)^2=5\ \Longrightarrow\ 5m^2=5\ \Longrightarrow\ m=\pm1 .\]

判别式 $\Delta=4m^2+4(m^2+2)=8m^2+8>0$，故两组均有两个不同交点。

\gkanswer{答案：（1）$x^2-\frac{y^2}{2}=1$；（2）$m=\pm1$。}

\end{answer}
\begin{answer}
\gksolhead{20}{}

\textbf{题目}：

设数列 $\{a_n\}$ 的通项公式为 $a_n=pn+q$（$n\in\mathbf N^*$，$P>0$）．数列 $\{b_n\}$ 定义如下：对于正整数 $m$，$b_m$ 是使得不等式 $a_n\geqslant m$ 成立的所有 $n$ 中的最小值．

（1）若 $p=\dfrac12$，$q=-\dfrac13$，求 $b_3$；

（2）若 $p=2$，$q=-1$，求数列 $\{b_m\}$ 的前 $2m$ 项和公式；

（3）是否存在 $p$ 和 $q$，使得 $b_m=3m+2$（$m\in\mathbf N^*$）？如果存在，求 $p$ 和 $q$ 的取值范围；如果不存在，请说明理由．

\Needspace{4\baselineskip}
\textbf{解析}

\Needspace{5\baselineskip}
（1）$a_n=\frac n2-\frac13$。由

\nopagebreak[4]
\[\frac n2-\frac13\geqslant3\ \Longrightarrow\ n\geqslant\frac{20}{3}\]

得最小整数 $n=7$，故 $b_3=7$。

\Needspace{5\baselineskip}
（2）$a_n=2n-1$。由 $2n-1\geqslant m$ 得 $n\geqslant\frac{m+1}{2}$，故

\nopagebreak[4]
\[b_m=\left\lceil\frac{m+1}{2}\right\rceil .\]

\Needspace{5\baselineskip}
于是

\nopagebreak[4]
\[b_{2k-1}=k,\qquad b_{2k}=k+1,\]

\Needspace{5\baselineskip}
前 $2m$ 项和

\nopagebreak[4]
\[S_{2m}=\sum_{k=1}^{m}(b_{2k-1}+b_{2k})
=\sum_{k=1}^{m}(2k+1)=m(m+1)+m=m^2+2m .\]

\Needspace{5\baselineskip}
（3）存在。$b_m=3m+2$ 等价于

\nopagebreak[4]
\[a_{3m+2}\geqslant m,\qquad a_{3m+1}<m .\]

\Needspace{5\baselineskip}
代入 $a_n=pn+q$：

\nopagebreak[4]
\[(3p-1)m+(2p+q)\geqslant0,\qquad
(3p-1)m+(p+q)<0 .\]

\Needspace{5\baselineskip}
若 $3p-1>0$，第二式左边随 $m$ 增大而增大，终将非负，矛盾；若 $3p-1<0$，第一式终将为负，矛盾。故必须

\nopagebreak[4]
\[p=\frac13 .\]

\Needspace{5\baselineskip}
此时两式化为

\nopagebreak[4]
\[\frac23+q\geqslant0,\qquad \frac13+q<0,\]

\Needspace{5\baselineskip}
即

\nopagebreak[4]
\[-\frac23\leqslant q<-\frac13 .\]

所以存在，且 $p=\frac13$，$q\in\left[-\frac23,-\frac13\right)$。

\gkanswer{答案：（1）$b_3=7$；（2）$m^2+2m$；（3）存在，$p=\frac13$，$q\in\left[-\frac23,-\frac13\right)$。}

\end{answer}